\chapter{Realisation et implementation}

\section{Introduction}
Ce chapitre decrit la mise en oeuvre technique de \ProjectName, depuis l'architecture logicielle jusqu'aux interfaces utilisateur exploitees par les equipes metier.

\section{Architecture technique de la solution}
La solution realisee s'appuie sur une architecture full-stack data orientee services.

\begin{table}[H]
\centering
\caption{Stack technique de \ProjectName}
\begin{tabular}{p{3.8cm}p{9.2cm}}
\toprule
\textbf{Couche} & \textbf{Technologies} \\
\midrule
Collecte & Python, Playwright, parseurs par site \\
Traitement & Python, Pandas, ETL hebdomadaire \\
Stockage & PostgreSQL \\
Services API & FastAPI, Pydantic, JWT \\
Interface & React, Vite, TypeScript, TanStack Query \\
\bottomrule
\end{tabular}
\end{table}

\figureplaceholder{Diagramme d'architecture technique complete}{Architecture technique de la plateforme}

\section{Collecte et transformation des donnees (Pipeline ETL)}
\subsection{Scraping automatise}
La collecte des prix est automatisee via des scrapers dedies aux sites cibles selon les priorites commerciales. Les donnees brutes sont stockees dans \texttt{raw\_staging} avec statut, code HTTP et capture d'ecran.

\subsection{Traitements ETL}
Le pipeline ETL nettoie les donnees, convertit les formats et consolide les prix dans les tables analytiques.

\begin{itemize}[leftmargin=1.2cm]
	\item validation des montants ;
	\item rejet des valeurs invalides (<= 0) ;
	\item normalisation des devises vers EUR ;
	\item calcul des agrégats hebdomadaires.
\end{itemize}

\begin{lstlisting}[language=Python,caption={Orchestration simplifiee du pipeline principal}]
if __name__ == "__main__":
    fetch_and_store_exchange_rates()
    run_all_scrapers()
    run_etl()
\end{lstlisting}

\figureplaceholder{Capture d'execution du pipeline en terminal}{Execution end-to-end du pipeline}

\section{Backend API et securite}
Le backend FastAPI expose les endpoints necessaires au dashboard : prix, filtres, series temporelles, administration et supervision operationnelle.

La securite est assuree par authentification token et controle d'acces base sur les roles :
\begin{itemize}[leftmargin=1.2cm]
	\item \textbf{admin} : acces complet y compris configuration ;
	\item \textbf{user} : consultation analytique.
\end{itemize}

\figureplaceholder{Capture API docs / endpoints principaux}{Services exposes par FastAPI}

\section{Presentation des interfaces}
\subsection{Vues analytiques}
\begin{itemize}[leftmargin=1.2cm]
	\item \textbf{Prices Explorer} : exploration detaillee des prix avec filtre et source URL.
	\item \textbf{Price Trends} : evolution hebdomadaire des prix normalises.
	\item \textbf{Store Comparison} : comparaison concurrentielle par enseigne/magasin.
\end{itemize}

\subsection{Vues administratives}
\begin{itemize}[leftmargin=1.2cm]
	\item gestion du catalogue produits (brands/categories/ranges/formats) ;
	\item gestion des URLs de collecte ;
	\item matrice Retail Presence ;
	\item supervision des echecs (Operational Visibility) ;
	\item gestion des utilisateurs.
\end{itemize}

\figureplaceholder{Capture dashboard principal}{Vue globale de l'interface}
\figureplaceholder{Capture page Retail Presence}{Pilotage de la couverture produit}
\figureplaceholder{Capture page Operational Visibility}{Supervision des incidents de collecte}

\section{Mise en oeuvre pratique}
L'exploitation de la plateforme suit cinq etapes :
\begin{enumerate}[leftmargin=1.2cm]
	\item preparation de l'environnement Python et installation des dependances ;
	\item initialisation de la base PostgreSQL ;
	\item chargement des donnees initiales ;
	\item execution du pipeline ;
	\item lancement backend et frontend.
\end{enumerate}

\section{Contributions des etudiants}
Le projet a ete realise en co-construction par \textbf{Khalil Amamri} et \textbf{Mahdi ElhadjAmor}, avec une implication sur la conception, l'implementation et la validation.

\section{Conclusion}
La phase de realisation confirme la faisabilite et la valeur operationnelle de \ProjectName. Le systeme relie des donnees web heterogenes a des indicateurs decisionnels directement exploitables par l'entreprise.

\clearpage
